%==============================================================================
%  DeltaV Interceptor -- Actuator & Aerodynamic Block Equations
%
%  Transcribed directly from the implementation:
%      src/actuators/rotors.py
%      src/actuators/lifting_body.py
%      src/actuators/elevons.py
%
%  Every equation below is the one the simulator actually evaluates, including
%  the regularisation and blending terms. Compile with pdflatex or xelatex.
%==============================================================================
\documentclass[11pt,a4paper]{article}

\usepackage[margin=2.4cm]{geometry}
\usepackage{amsmath,amssymb,bm}
\usepackage{booktabs}
\usepackage{longtable}
\usepackage{array}
% xcolor must be loaded ONCE with every option it needs -- loading it twice with
% different options is an option clash, and hyperref cannot reference a colour
% that does not exist yet, so xcolor and the colour definition both come first.
\usepackage[table,usenames,dvipsnames]{xcolor}
\definecolor{blockrule}{RGB}{29,93,135}

\usepackage{titlesec}
\usepackage{enumitem}
\usepackage[colorlinks=true,linkcolor=blockrule,citecolor=blockrule,
            urlcolor=blockrule]{hyperref}

\titleformat{\section}{\Large\bfseries\color{blockrule}}{\thesection}{0.7em}{}
\titleformat{\subsection}{\large\bfseries}{\thesubsection}{0.6em}{}

\newcommand{\vect}[1]{\bm{#1}}
\newcommand{\unit}[1]{\hat{\bm{#1}}}
\newcommand{\dd}{\,\mathrm{d}}
\newcommand{\Rey}{\mathrm{Re}}
\newcommand{\clip}{\operatorname{clip}}
\newcommand{\sgn}{\operatorname{sgn}}
\newcommand{\atantwo}{\operatorname{atan2}}
\newcommand{\interp}{\operatorname{interp}}

\title{\textbf{DeltaV Interceptor}\\[2mm]
\large Governing Equations of the Rotor, Lifting-Body and Elevon Blocks}
\author{6\,DOF Actuator Testbench}
\date{}

\begin{document}
\maketitle

\begin{abstract}
\noindent
This document sets out the complete set of equations evaluated by the three
aerodynamic/propulsive blocks of the DeltaV Interceptor 6\,DOF model. The
vehicle is a quad tilt-rotor flying wing: four independently tilting rotors on
an asymmetric H-frame, a swept tailless wing with two elevon channels, and no
separate control surfaces beyond those. Sign conventions follow the body frame
$x$ forward, $y$ right, $z$ down (NED). All equations are transcribed from the
implementation, including the regularisation, blending and validity-fade terms
that make the model well behaved at hover, through stall, and beyond the
tabulated data envelope.
\end{abstract}

%==============================================================================
\section*{Nomenclature}
\addcontentsline{toc}{section}{Nomenclature}

\begin{longtable}{@{}p{2.6cm}p{6.1cm}p{2.4cm}p{4.6cm}@{}}
\toprule
\textbf{Symbol} & \textbf{Meaning} & \textbf{Symbol} & \textbf{Meaning}\\
\midrule
\endhead
$u,v,w$            & body-frame CG velocity            & $p,q,r$          & body angular rates \\
$V_a$              & true airspeed                     & $\bar q$         & freestream dynamic pressure \\
$\alpha,\beta$     & body incidence, sideslip          & $\alpha_w$       & wing incidence, post-downwash \\
$\varepsilon$      & rotor downwash angle              & $\bar q_w$       & dynamic pressure at the wing \\
$\rho$             & air density                       & $S,b,\bar c$     & wing area, span, mean chord \\
$A\!R$             & aspect ratio                      & $e$              & Oswald efficiency \\
$n_i$              & rotor speed, rev\,s$^{-1}$        & $\lambda_i$      & rotor tilt angle \\
$D$                & rotor diameter                    & $R$              & rotor radius, $D/2$ \\
$A$                & rotor disc area                   & $s_i$            & spin sign, $\pm1$ \\
$x_i,y_i$          & rotor station about the CG        & $J_i$            & advance ratio \\
$C_T,C_Q$          & thrust, torque coefficients       & $T_i,Q_i$        & rotor thrust, torque \\
$\mu_i$            & in-plane advance ratio            & $\alpha_i$       & disc incidence \\
$I_p$              & rotor polar inertia               & $\sigma_b,C_{d}$ & blade solidity, blade drag \\
$\delta_1,\delta_2$& elevon channel deflections        & $\delta_e,\delta_a$ & pitch, roll elevon inputs \\
$\eta$             & DATCOM deflection efficiency      & $y_{\!e}$        & elevon spanwise centroid \\
$\alpha_{s}$       & stall angle                       & $\sigma$         & stall blend fraction \\
$x_{np},x_{cg}$    & neutral point, CG station         & $k_w$            & wing activity fraction \\
$V_{\!\mathrm{reg}}$ & airspeed regularisation floor   & $n_{\mathrm{reg}}$ & rotor-speed floor \\
\bottomrule
\end{longtable}

\noindent
\textbf{Regularisation constants.} $V_{\!\mathrm{reg}}=5\ \mathrm{m\,s^{-1}}$,
$V_{db}=0.5\ \mathrm{m\,s^{-1}}$ (hover deadband),
$\epsilon_u=1\ \mathrm{m\,s^{-1}}$, $n_{\mathrm{reg}}=5\ \mathrm{rev\,s^{-1}}$,
$\varepsilon_{\max}=85^\circ$, $\Delta_{ps}=2^\circ$ (post-stall blend offset).

%==============================================================================
\section{Rotor Block}
\emph{Source:} \texttt{src/actuators/rotors.py}

\medskip\noindent
Four rotors, indexed $i=1\ldots4$ in the order
$[\,\mathrm{FR},\ \mathrm{AL},\ \mathrm{FL},\ \mathrm{AR}\,]$
(front-right, aft-left, front-left, aft-right), with spin signs
$s_i=[+1,+1,-1,-1]$. The tilt convention is
$\lambda=0$ for hover (thrust along $-z_b$) and $\lambda=\pi/2$ for cruise
(thrust along $+x_b$), giving the shaft unit vector

\begin{equation}
\unit n_i =
\begin{bmatrix}\sin\lambda_i\\[1mm] 0 \\[1mm] -\cos\lambda_i\end{bmatrix}.
\end{equation}

\subsection{Hub kinematics}
Each hub is offset from the CG, so it sees the CG velocity plus the rigid-body
rotation carrying it around, $\vect V_{h,i}=\vect V_{cg}+\vect\Omega\times\vect r_i$
with $\vect r_i=(x_i,y_i,0)$:

\begin{equation}
\vect V_{h,i}=
\begin{bmatrix}
u - r\,y_i \\
v + r\,x_i \\
w + p\,y_i - q\,x_i
\end{bmatrix}.
\end{equation}

The axial inflow is its projection onto the shaft, $V_{ax,i}=\vect V_{h,i}\cdot\unit n_i$.
Since $\unit n_i$ has no $y$-component (tilt is single-axis about body $y$),

\begin{equation}
\boxed{\;
V_{ax,i}=\bigl(u-r\,y_i\bigr)\sin\lambda_i-\bigl(w+p\,y_i-q\,x_i\bigr)\cos\lambda_i \;}
\end{equation}

\subsection{Advance ratio and coefficient lookup}
\begin{equation}
J_i=\frac{V_{ax,i}}{\max\!\left(n_i,\,n_{\mathrm{reg}}\right)D}
\end{equation}

$C_T$ and $C_Q$ are tabulated against $J$. Below the table the lookup clamps
(correct: static thrust). \emph{Above} the last tabulated point $J_N$ the table
is continued linearly into the windmill-brake regime rather than clamped, so a
propeller past its zero-thrust advance ratio correctly produces negative thrust:

\begin{equation}
C(J)=
\begin{cases}
\interp\!\left(J;\,J_k,\,C_k\right), & J\le J_N\\[3mm]
\max\!\Bigl(C_N+m\,\min\!\left(J-J_N,\ \Delta J_{\max}\right),\ C_{\mathrm{floor}}\Bigr), & J>J_N
\end{cases}
\qquad
m=\frac{C_N-C_{N-1}}{J_N-J_{N-1}}
\end{equation}

with $\Delta J_{\max}=1.0$ and $C_{\mathrm{floor}}=-0.06$.

\subsection{Thrust, torque and their moments}
\begin{align}
T_i&=C_T(J_i)\,\rho\,n_i^{2}D^{4}, &
Q_i&=C_Q(J_i)\,\rho\,n_i^{2}D^{5}
\end{align}

Resolved into the body frame, $\vect F_{T,i}=T_i\unit n_i$:
\begin{align}
F_{Tx,i}&=T_i\sin\lambda_i, &
F_{Tz,i}&=-T_i\cos\lambda_i
\end{align}

The thrust moments follow from $\vect r_i\times\vect F_{T,i}$:
\begin{align}
M_{Tx,i}&=-y_i\,T_i\cos\lambda_i &&\text{(roll)}\\
M_{Ty,i}&=+x_i\,T_i\cos\lambda_i &&\text{(pitch)}\\
M_{Tz,i}&=-y_i\,T_i\sin\lambda_i &&\text{(yaw)}
\end{align}

The motor reaction torque acts along $-s_i\,Q_i\unit n_i$:
\begin{align}
M_{Qx,i}&=-s_i\,Q_i\sin\lambda_i, &
M_{Qz,i}&=+s_i\,Q_i\cos\lambda_i
\end{align}

\subsection{Oblique-flow correction (normal force and P-factor)}
The axial table cannot see the in-plane component of the hub velocity. Define

\begin{align}
\vect V_{\perp,i}&=\vect V_{h,i}-V_{ax,i}\,\unit n_i, &
\unit d_i&=\frac{\vect V_{\perp,i}}{\lVert\vect V_{\perp,i}\rVert}
\end{align}

\begin{equation}
\alpha_i=
\begin{cases}
0, & \lVert\vect V_{\perp,i}\rVert<10^{-9}\\[1mm]
\atantwo\!\left(\lVert\vect V_{\perp,i}\rVert,\ V_{ax,i}\right), & \text{otherwise}
\end{cases}
\qquad
\mu_i=\frac{\lVert\vect V_{\perp,i}\rVert}{\Omega_i R},
\quad \Omega_i=2\pi\max(n_i,n_{\mathrm{reg}})
\end{equation}

The disc dynamic pressure and the singular-term guard are
\begin{align}
q_i&=\tfrac12\rho\,\max\!\left(\lVert\vect V_{h,i}\rVert,\ V_{\!\mathrm{reg}}\right)^{2}, &
J_{s}&=\sgn(J_i)\max\!\left(|J_i|,\ J_{\mathrm{reg}}\right)
\end{align}

The mean blade lift coefficient is
\begin{equation}
\bar C_{l,i}=\frac{3J_i}{2\pi}
\left[\frac{2}{\sigma_b q_i A}\cdot\frac{J_i}{\pi}\,T_i+C_{d}\right]
\end{equation}

and the normal force and P-factor moment magnitudes, with $a=2\pi$ the
thin-airfoil lift-curve slope, are
\begin{align}
P_{N,i}&=\frac{\sigma_b q_i A}{2}
\left[\bar C_{l,i}
+\frac{a J_i}{2\pi}\ln\!\left(1+\Bigl(\frac{\pi}{J_{s}}\Bigr)^{2}\right)
+\frac{\pi}{J_{s}}C_{d}\right]\alpha_i\,\Lambda_i\\[2mm]
N_{P,i}&=-\frac{\sigma_b q_i A R}{2}
\left[\frac{2\pi}{3J_{s}}\bar C_{l,i}
+\frac{a}{2}\left(1-\Bigl(\frac{J_{s}}{\pi}\Bigr)^{2}
\ln\!\left(1+\Bigl(\frac{\pi}{J_{s}}\Bigr)^{2}\right)\right)
-\frac{\pi}{J_{s}}C_{d}\right]\alpha_i\,\Lambda_i
\end{align}

Both are gated by a combined validity factor $\Lambda_i$ built from a
magnitude ramp (so a vanishing $\vect V_\perp$ cannot hand the correction a
numerically arbitrary direction, since $\alpha_i$ is an \emph{angle} and can sit
at $90^\circ$ for an arbitrarily small $\vect V_\perp$) and an in-plane advance
ratio fade (the linear model is trusted only below $\mu\approx0.3$):

\begin{equation}
\Lambda_i=
\underbrace{\clip\!\left(\frac{\lVert\vect V_{\perp,i}\rVert}{V_{\!\mathrm{reg}}},0,1\right)}_{\text{magnitude ramp}}
\cdot
\underbrace{\clip\!\left(1-\frac{\mu_i-0.3}{0.4-0.3},\,0,\,1\right)}_{\text{validity fade}}
\end{equation}

\begin{align}
\vect F_{N,i}&=P_{N,i}\,\unit d_i, &
\vect M_{N,i}&=s_i\,N_{P,i}\,\unit d_i
\end{align}

\subsection{Gyroscopic coupling}
Each shaft carries angular momentum $\vect h_i=I_p\,(2\pi n_i)\,s_i\,\unit n_i$.
With the shaft angular velocity including the tilt rate,
$\vect\omega_{s,i}=\vect\omega_{b}+\dot\lambda_i\,\unit e_{\mathrm{tilt}}$,

\begin{equation}
\vect M_{\mathrm{gyro}}=-\sum_{i=1}^{4}\vect\omega_{s,i}\times\vect h_i
\end{equation}

\subsection{Block totals}
\begin{align}
\vect F_{\mathrm{rot}}&=\sum_{i=1}^{4}
\begin{bmatrix}F_{Tx,i}+F_{Nx,i}\\ F_{Ny,i}\\ F_{Tz,i}+F_{Nz,i}\end{bmatrix}\\[2mm]
\vect M_{\mathrm{rot}}&=\sum_{i=1}^{4}
\begin{bmatrix}
M_{Tx,i}+M_{Nx,i}+M_{Qx,i}\\
M_{Ty,i}+M_{Ny,i}\\
M_{Tz,i}+M_{Nz,i}+M_{Qz,i}
\end{bmatrix}+\vect M_{\mathrm{gyro}}
\end{align}

The front-pair thrust handed to the wing block (Sec.~\ref{sec:downwash}) is
\begin{equation}
T_f=T_1+T_3
\end{equation}

%==============================================================================
\section{Lifting-Body Block}
\emph{Source:} \texttt{src/actuators/lifting\_body.py}

\subsection{Air data}
\begin{align}
V_a&=\sqrt{u^{2}+v^{2}+w^{2}}, &
\beta&=\arcsin\!\left(\clip\!\left(\frac{v}{V_a},-1,1\right)\right), &
\bar q&=\tfrac12\rho V_a^{2}
\end{align}

The incidence $\alpha$ is regularised, because $\atantwo(w,u)$ is genuinely
singular as $V_a\to0$ (the direction of a near-zero vector is not physically
meaningful) and has a branch cut at $w=0,\ u<0$. Two devices are used: a
softplus floor on $u$, and a linear blend toward a bounded reference:

\begin{align}
u_{f}&=\tfrac12\!\left(u+\sqrt{u^{2}+\epsilon_u^{2}}\right)
&&\text{smooth }\max(u,0),\ \text{never exactly }0\\
\alpha_{\mathrm{raw}}&=\atantwo\!\left(w,\ u_f\right), &
\alpha_{\mathrm{reg}}&=\atantwo\!\left(w,\ V_{\!\mathrm{reg}}\right)\\
k&=\min\!\left(V_a/V_{\!\mathrm{reg}},\,1\right), &
\alpha&=(1-k)\,\alpha_{\mathrm{reg}}+k\,\alpha_{\mathrm{raw}}
\end{align}

\subsection{Wing activity}
The wing's entire contribution is faded out at negligible airspeed, with a hard
floor below the deadband so no self-feeding drift can survive at true hover:
\begin{equation}
k_w(V_a)=
\begin{cases}
0, & V_a<V_{db}\\
\min\!\left(V_a/V_{\!\mathrm{reg}},\,1\right), & \text{otherwise}
\end{cases}
\end{equation}

\subsection{Rotor downwash}\label{sec:downwash}
Only the component of front-rotor thrust still pointing \emph{through} the wing
contributes; the $\cos\lambda_f$ gate fades this to zero at cruise tilt:
\begin{align}
T_{f\perp}&=\max\!\left(T_f\cos\lambda_f,\ 0\right), &
w_i&=\sqrt{\frac{T_{f\perp}}{2\rho A}} &&\text{(momentum theory)}
\end{align}

\begin{equation}
\varepsilon=
\begin{cases}
0, & V_a<V_{db}\\[1mm]
\clip\!\left(\dfrac{k_\varepsilon\,w_i}{\max(V_a,V_{\!\mathrm{reg}})},
\ -\varepsilon_{\max},\ \varepsilon_{\max}\right), & \text{otherwise}
\end{cases}
\end{equation}

\begin{equation}
\bar q_w=
\begin{cases}
\bar q\left(1+\dfrac{k_q\,T_{f\perp}}{\bar q\,A}\right), & \bar q>10^{-9}\\[3mm]
\dfrac{k_q\,T_{f\perp}}{A}, & \text{hover, no freestream}
\end{cases}
\qquad\Longrightarrow\qquad
\boxed{\;\alpha_w=\alpha-\varepsilon\;}
\end{equation}

\subsection{Lift and drag: attached, separated, and the blend}
The stall blend is a smooth two-sided logistic, $0$ pre-stall and $1$ post-stall:
\begin{equation}
\sigma(\alpha_w)=\frac{1+e^{z_1}+e^{z_2}}{\left(1+e^{z_1}\right)\left(1+e^{z_2}\right)},
\qquad
z_1=-M\!\left(\alpha_w-\alpha_s^{*}\right),\quad
z_2=+M\!\left(\alpha_w+\alpha_s^{*}\right)
\end{equation}
with $\alpha_s^{*}=\alpha_{s}+\Delta_{ps}$. The offset places the \emph{peak} of
the resulting lift curve on $\alpha_s$; a blend centred on $\alpha_s$ would put
$\sigma=\tfrac12$ at the peak and roughly halve $C_{L,\max}$.

\paragraph{Attached branch.} Measured polar data where it is trusted, parametric
model outside it. With table range $[\alpha_{\min},\alpha_{\max}]$ and taper
width $\Delta_b$, the trust weight is a cosine ease:
\begin{equation}
W=
\begin{cases}
1, & \alpha_{\min}\le\alpha_w\le\alpha_{\max}\\[1mm]
0, & d\ge\Delta_b\\[1mm]
\tfrac12-\tfrac12\cos\!\left(\pi\left(1-d/\Delta_b\right)\right), & \text{otherwise}
\end{cases}
\end{equation}
where $d$ is the distance past whichever edge was crossed. The parametric
partner is
\begin{align}
C_{L}^{\,\mathrm{par}}&=C_{L0}+C_{L\alpha}\,\alpha_w, &
C_{D}^{\,\mathrm{par}}&=C_{Dp}+\frac{\bigl(C_{L}^{\,\mathrm{par}}\bigr)^{2}}{\pi e\,A\!R}
\end{align}
and the attached values are $C^{\,\mathrm{att}}=W\,C^{\,\mathrm{tab}}+(1-W)\,C^{\,\mathrm{par}}$.

\paragraph{Separated branch.} A stalled wing behaves as a flat plate carrying a
force normal to its surface, with the finite-aspect-ratio normal-force
coefficient (Hoerner correlation)
\begin{equation}
C_N=1.11+0.018\,A\!R
\end{equation}
\begin{align}
C_{L}^{\,\mathrm{sep}}&=C_N\sin\alpha_w\cos\alpha_w, &
C_{D}^{\,\mathrm{sep}}&=C_{Dp}+C_N\sin^{2}\alpha_w
\end{align}
so $C_L\to0$ and $C_D\to C_N$ at $\alpha_w=90^\circ$, both correct.

\paragraph{Result.} Stall is applied once, identically for either data source,
and the non-lifting airframe drag is added on top:
\begin{align}
C_L&=(1-\sigma)\,C_{L}^{\,\mathrm{att}}+\sigma\,C_{L}^{\,\mathrm{sep}}\\
C_D&=(1-\sigma)\,C_{D}^{\,\mathrm{att}}+\sigma\,C_{D}^{\,\mathrm{sep}}
+\underbrace{C_{D,\mathrm{airframe}}}_{\text{fuselage, nacelles, booms}}
\end{align}

\subsection{Moment coefficients}
With $V_s=\max(V_a,V_{\!\mathrm{reg}})$,
\begin{align}
C_m&=C_{m0}+C_{m\alpha}\alpha_w+C_{mq}\frac{\bar c}{2V_s}\,q\\
C_l&=C_{l\beta}\,\beta+C_{lp}\frac{b}{2V_s}\,p\\
C_n&=C_{n\beta}\,\beta+C_{nr}\frac{b}{2V_s}\,r
\end{align}

\subsection{Forces and moments}
\begin{align}
L&=k_w\,\bar q_w S\left(C_L+C_{Lq}\frac{\bar c}{2V_s}q\right), &
D&=k_w\,\bar q_w S\,C_D, &
Y&=k_w\,\bar q_w S\,C_{Y\beta}\,\beta
\end{align}

The pitching moment is taken about the CG as
$\vect r_{cg\to np}\times\vect F$. With stations measured aft from the nose the
neutral point lies at $-(x_{np}-x_{cg})$ along body $+x$, leaving
$M_y=(x_{np}-x_{cg})F_z$, where $F_z$ is the \emph{body-axis normal force}, not
the wind-axis lift:
\begin{equation}
F_{z,w}=-\left(L\cos\alpha_w+D\sin\alpha_w\cos\beta+Y\sin\alpha_w\sin\beta\right)
\end{equation}
\begin{align}
M_x&=k_w\,\bar q_w S b\,C_l\\
M_y&=\underbrace{k_w\,\bar q_w S\bar c\,C_{m0}}_{\text{zero-lift couple}}
+\underbrace{\left(x_{np}-x_{cg}\right)F_{z,w}}_{\text{NP lever arm}}
+\underbrace{k_w\,\bar q_w S\bar c\,C_{mq}\frac{\bar c}{2V_s}q}_{\text{pitch damping}}\\
M_z&=k_w\,\bar q_w S b\,C_n
\end{align}

\noindent\textbf{Note.} $C_{m\alpha}$ is computed but deliberately \emph{not}
added to $M_y$: the neutral-point lever arm already \emph{is} the
$\alpha$-dependent pitching moment, and including both would double-count it.
The same reasoning applies to yaw, where $C_{n\beta}$ carries the whole
weathercock effect and no separate sideforce-times-arm term is added.

\subsection{Projection to body axes}
The wind-axis forces are rotated through $\alpha_w$ and $\beta$:
\begin{align}
F_x&=L\sin\alpha_w-D\cos\alpha_w\cos\beta-Y\cos\alpha_w\sin\beta\\
F_y&=Y\cos\beta-D\sin\beta\\
F_z&=-L\cos\alpha_w-D\sin\alpha_w\cos\beta-Y\sin\alpha_w\sin\beta
\end{align}

The moments are passed through \emph{unrotated}: $C_l$, $C_m$, $C_n$ and the
lever-arm term are body-axis quantities by definition ($C_{lp}$ multiplies body
roll rate $p$, $C_{nr}$ multiplies body yaw rate $r$), so applying the
wind-to-body rotation a second time would be incorrect.

%==============================================================================
\section{Elevon Block}
\emph{Source:} \texttt{src/actuators/elevons.py}

\subsection{Channel split}
The two physical surfaces are decomposed into a symmetric (pitch) and a
differential (roll) input:
\begin{align}
\delta_e&=\tfrac12\left(\delta_1+\delta_2\right) &&\text{pitch}\\
\delta_a&=\tfrac12\left(\delta_2-\delta_1\right) &&\text{roll}
\end{align}

\subsection{Effectiveness}
Large-deflection efficiency is interpolated from a DATCOM table,
$\eta(\delta)=\interp\!\left(|\delta|;\ \eta_{\mathrm{tab}}\right)$.
This is then knocked down by the wing's own separation state, using the
\emph{same} $\sigma$ as the lifting-body block: a trailing-edge surface only
works while the flow ahead of it is attached.
\begin{equation}
k_{\mathrm{auth}}=1-\left(1-\eta_{ps}\right)\sigma(\alpha_w),
\qquad \eta_{ps}=0.15
\end{equation}
\begin{align}
\eta_e&\leftarrow\eta(\delta_e)\,k_{\mathrm{auth}}, &
\eta_a&\leftarrow\eta(\delta_a)\,k_{\mathrm{auth}}
\end{align}

A hard operating limit is applied separately: at or above the stall cap,
trailing-edge-down deflection is forbidden, since it would push the already
near-separated aft wing further past stall,
\begin{equation}
\delta_e\leftarrow\max\left(\delta_e,\,0\right)
\qquad\text{if}\qquad \alpha_w\ge\alpha_{s}
\end{equation}

\subsection{Incremental forces and moments}
Evaluated in the \emph{slipstream-corrected} $\bar q_w$ and $\alpha_w$, since the
elevons sit in the rotor wash rather than in the freestream:
\begin{align}
\Delta L&=-C_{L\delta_e}\,\eta_e\,\delta_e\,\bar q_w S
&&\text{lift loss from deflection}\\
\Delta D_e&=\bar q_w S\,k_e\left(\eta_e\delta_e\right)^{2}
&&\text{airbrake drag}
\end{align}
\begin{align}
F_x&=\Delta L\sin\alpha_w-\Delta D_e\cos\alpha_w\\
F_y&=0\\
F_z&=-\Delta L\cos\alpha_w-\Delta D_e\sin\alpha_w
\end{align}
\begin{align}
M_x&=2\,\bar q_w S\,C_{L\delta_e}\,\eta_a\,\delta_a\,y_{\!e} &&\text{roll}\\
M_y&=C_{m\delta_e}\,\eta_e\,\delta_e\,\bar q_w S\,\bar c &&\text{pitch}\\
M_z&=0 &&\text{(finless airframe, no rudder)}
\end{align}

%==============================================================================
\section{Assembly}
The three blocks are summed once, in body axes, and passed to the rigid-body
equations. Gravity enters the force sum only, never the moment sum, which is
equivalent to assuming the reference point coincides with the CG:
\begin{align}
\sum\vect F&=\vect F_{\mathrm{rot}}+\vect F_{\mathrm{wing}}+\vect F_{\mathrm{elev}}\\
\sum\vect M&=\vect M_{\mathrm{rot}}+\vect M_{\mathrm{wing}}+\vect M_{\mathrm{elev}}
\end{align}
\begin{align}
\dot{\vect V}_b&=\frac{1}{m}\left(\mathbf{R}^{\!\top}m\vect g
-m\,\vect\omega\times\vect V_b+\sum\vect F\right)\\
\dot{\vect\omega}&=\mathbf{J}^{-1}\left(\sum\vect M-\vect\omega\times\mathbf{J}\vect\omega\right),
\qquad
\mathbf{J}=\begin{bmatrix}
I_{xx}&0&-I_{xz}\\ 0&I_{yy}&0\\ -I_{xz}&0&I_{zz}
\end{bmatrix}\\
\dot{\vect\Lambda}&=\tfrac12\,\Omega(\vect\omega)\,\vect\Lambda
\qquad\text{(quaternion kinematics)}
\end{align}

\noindent
The coupling that prevents any block being evaluated independently is the chain
\begin{equation*}
T_f \;\longrightarrow\; \varepsilon,\ \bar q_w
\;\longrightarrow\; \alpha_w \;\longrightarrow\; C_L
\;\longrightarrow\; M_y \;\longrightarrow\; \text{attitude}
\;\longrightarrow\; \vect V_b \;\longrightarrow\; V_{ax,i},\ \alpha,
\end{equation*}
which closes through the rotor inflow and the wing incidence simultaneously.

\end{document}
